\documentclass[12pt]{article}
\usepackage[T1]{fontenc}
\usepackage[utf8]{inputenc}
\usepackage{lmodern}
\usepackage{textcomp}
\usepackage{enumitem}
\usepackage{geometry}
\geometry{margin=1in}
\begin{document}
\begin{center}
Answer Key - Philosophy - Graduate Level Assessment 15 \
\small (Answers are ordered exactly as in the question paper.)
\end{center}

\section*{Multiple Choice}
\begin{enumerate}[label=\arabic*.]
\item \textbf{Answer:} A
\\ \textit{Options:} A) a valid modus ponens; B) a valid modus tollens; C) a valid dilemma; D) a valid disjunctive syllogism; E) a valid conditional syllogism
\\ \textit{Note:} The requirement that "each term must be used twice, no more and no less" reflects the structure of valid modus ponens, where the conditional statement and its antecedent are each applied once to derive the consequent, ensuring a balanced logical progression.
\item \textbf{Answer:} A
\\ \textit{Options:} A) provides some evidence against psychological egoism.; B) supports the theory of psychological altruism.; C) is irrelevant to the discussion of psychological egoism.; D) supports the idea of ethical egoism.; E) proves the concept of psychological egoism.
\\ \textit{Note:} Feinberg argues that Lincoln's selfless act of sacrificing his own interests for the benefit of others suggests that individuals can act out of genuine concern for others, thereby providing evidence against the notion of psychological egoism, which posits that all actions are ultimately motivated by self-interest.
\item \textbf{Answer:} E
\\ \textit{Options:} A) Tradition; B) Clerical authority; C) Ethical principles; D) Historical context; E) Divine inspiration
\\ \textit{Note:} Divine inspiration is not included in the Islamic understanding of jurisprudence because Islamic law, or Sharia, is derived from the Quran, Hadith, consensus (ijma), and analogy (qiyas), rather than direct divine inspiration to human jurists.
\item \textbf{Answer:} A
\\ \textit{Options:} A) Wealth and prosperity; B) Fear and obedience; C) Knowledge and education; D) Power and influence; E) Truth and love
\\ \textit{Note:} Wealth and prosperity have been central focuses of religious traditions in China and Korea due to the Confucian emphasis on social harmony, moral development, and the role of material success in achieving a stable and prosperous society.
\item \textbf{Answer:} A
\\ \textit{Options:} A) how "justice" is to be defined.; B) whether ethics is a science.; C) whether morality is subjective.; D) whether God exists.; E) whether humans are inherently good or bad.
\\ \textit{Note:} Moore argues that defining "justice" is central to ethical discussions because it serves as a foundational concept that influences the interpretation and application of moral principles across various contexts.
\end{enumerate}

\section*{True / False}
\begin{enumerate}[label=\arabic*.]
\setcounter{enumi}{5}
\item \textbf{Answer:} False
\\ \textit{Options:} A) True; B) False
\\ \textit{Note:} Cicero emphasizes that the most practical branch of philosophy is ethics, which deals with moral conduct and the good life, rather than epistemology, which focuses on knowledge and belief.
\item \textbf{Answer:} False
\\ \textit{Options:} A) True; B) False
\\ \textit{Note:} Kant asserts that laws of nature govern phenomena and causal relations in the physical world, while laws of freedom pertain to the moral law guiding rational agents, meaning that laws of nature do not imply the absence of events but rather the framework within which they occur.
\item \textbf{Answer:} True
\\ \textit{Options:} A) True; B) False
\\ \textit{Note:} Gauthier's moral framework posits that justified constraint is critical for establishing cooperative agreements among rational agents, thereby providing a foundation for resolving ethical dilemmas and addressing crises in moral philosophy.
\item \textbf{Answer:} False
\\ \textit{Options:} A) True; B) False
\\ \textit{Note:} The statement is false because the synoptic Gospels specifically refer to Matthew, Mark, and Luke, while John is considered a distinct Gospel with a unique theological perspective.
\item \textbf{Answer:} True
\\ \textit{Options:} A) True; B) False
\\ \textit{Note:} Wellman contends that without collective enforcement mechanisms, the absolute assertion of unlimited property rights undermines social order by allowing individuals to act solely in their self-interest, which can lead to a breakdown of cooperation and the emergence of an anarchic state.
\end{enumerate}

\section*{Long Form Response}
\begin{enumerate}[label=\arabic*.]
\setcounter{enumi}{10}
\item \textbf{Answer:} The fallacy of ignorance of refutation occurs when an individual either genuinely lacks the ability to engage in effective refutation or deliberately feigns this inability, leading to significant confusion in argumentative discourse. This lack of engagement can obscure the logical structure of arguments, allowing unsupported claims to proliferate unchallenged, thus undermining critical dialogue. When participants in a debate fail to adequately address counterarguments, it creates an environment where misconceptions flourish and the discourse devolves into mere assertion rather than reasoned argumentation. Ultimately, this fallacy not only hampers the clarity of discussions but also stifles intellectual growth, as the inability or refusal to refute effectively prevents the rigorous examination of ideas essential for philosophical advancement.
\\ \textit{Note:} correct because it accurately identifies that the fallacy of ignorance of refutation hinders effective argumentation by allowing unchallenged claims to persist, thereby obscuring logical discourse and preventing meaningful engagement with counterarguments.
\item \textbf{Answer:} Ziran, often translated as "naturalness" or "spontaneity," is a central concept in Zhuangzi's philosophy that emphasizes the inherent, uncontrived nature of existence. In this framework, compassion emerges as a natural response when individuals align themselves with the flow of the Dao, recognizing the interconnectedness of all beings. Zhuangzi suggests that when one acts from ziran, compassion is not a moral obligation but an organic expression of one's true self, free from artificial constructs of societal norms. This perspective invites a profound understanding of compassion as an intrinsic quality of being, fostering a sense of empathy that transcends conventional ethical considerations. Ultimately, the implications of ziran challenge rigid moral frameworks, advocating for a compassionate engagement with the world that is both spontaneous and harmonious.
\\ \textit{Note:} correct because it accurately captures how ziran, as spontaneity and naturalness, allows for compassion to arise organically from an individual's authentic existence and recognition of interconnectedness, rather than as a forced moral duty, aligning with Zhuangzi's philosophical emphasis on harmony with the Dao.
\end{enumerate}

\vfill
\noindent\textit{End of Answer Key}
\end{document}